% !TITLE 元
% !SUMMARY 俗文学的崛起
% !ORDER 0

\begin{intro}
元代不足百年，却是中国文学从雅向俗深刻转向的时代。

蒙古人入主中原，科举长期停废，传统文人的上升通道被阻断。一大批知识分子流落市井，与勾栏瓦舍的艺人混居，反而催生了一种全新的文学形态——元曲。元曲包括散曲和杂剧，这里要说的是散曲：它吸收了民间小调、少数民族音乐的元素，语言直白泼辣，风格诙谐洒脱，与诗词的含蓄典雅形成了鲜明的对照。

“枯藤老树昏鸦，小桥流水人家，古道西风瘦马。夕阳西下，断肠人在天涯。”马致远的《天净沙·秋思》是元曲小令的巅峰之作。全曲无一动词，只是名词的叠加，却像电影镜头一样推进，在最后一句点出一个漂泊者的孤独。这种“列锦”的手法，是诗词中罕见的。

张养浩的《山坡羊·潼关怀古》则展现了元曲的另一面：深沉的历史感。“兴，百姓苦；亡，百姓苦”——这不是文人的无病呻吟，而是一个历经宦海沉浮的老者对历史的透彻洞察。元曲里有很多这样的句子，粗看直白如话，细品却力透纸背。

元曲的价值，在于它打破了文学的精英垄断。在它之前，诗歌是士大夫的专利；在它之后，文学开始向更广阔的社会阶层开放。这种世俗化、平民化的倾向，直接影响了明清小说和戏曲的发展。
\end{intro}
